\section{Discussion}\label{sec:discussion}

\paragraph{What the evidence supports.} Time-uniform false-revocation control held in every setting we tried, including a planner that watches the monitor and a stream whose state distribution is a function of the model. A first-order delay law predicts measured delay across horizons to within $20\%$ on a system where the post-change statistic is known (the finite-sample bound of Theorem~\ref{thm:delay} is $9$--$15\times$ loose there). The certified total-variation bound was never violated and was off by a factor of $1.6$ to $3$ after a change. On
real video the construction detects every shift with a zero margin, at false-alarm rates of $0.43$--$0.50$ for the linear model; margins that bring false alarms to about $0.03$ cost detection on the linear model ($0.33$--$0.70$), but a frozen-encoder neural model reaches zero false alarms with detection $0.90$--$1.00$ at the same margin, and a weak convolutional autoencoder cannot be repaired by any margin. With a neural dynamics model in a closed loop, a tolerance-based gate with a moderate margin removed most of the catastrophic cost of an unfaithful long-horizon model with no false trigger in $100$ fault-free episodes, but it was slower than a CUSUM gate on the abrupt fault, and its advantage on the other two faults comes from not reacting to changes that the fallback would make worse. The tolerance and the quality of the representation decide the trade-off.

\paragraph{Limitations.}
\begin{itemize}[leftmargin=1.4em,itemsep=2pt]
\item \emph{The tolerance is a modeling choice and the dangerous one.} The guarantee needs $\E[Z_k\mid\Fcal_{k-1}]\le\eps_h$ conditionally. A constant tolerance set to the pooled in-distribution mean of a statistic whose conditional mean varies with the state does not satisfy it (ablation D: $0.253$ false revocations; real video with $k=0$: Section~\ref{sec:exp-real}). The standardized Gaussian excess is pivotal and has no such problem, but only for models that output a calibrated predictive covariance.
On real video we calibrate $\eps_h$ from held-out clips with a margin; that is empirically safe in our replications and not a guarantee.
\item \emph{Certificates are expensive.} Pinsker's inequality squares twice on the way from the mean of the statistic to total variation, so certifying a budget $\tau$ needs on the order of $\tau^{-4}$ rounds at one horizon and $h$ times more at horizon $h$.
In the simulator, windows of a few hundred rounds certify a gap of the size of our drifts but not a small budget on a faithful model (certified $1.63$ nats against a true $0$ in Figure~\ref{fig:teaser}). The certified horizon is therefore informative after large changes and conservative otherwise.
\item \emph{Long horizons are intrinsically costly.} Non-overlapping chains give horizon $h$ a $1/h$ share of the rounds; we found no valid way to pool the discarded rounds (Section~\ref{sec:overlap}). In our simulator a growing margin offsets this; in general it need not.
\item \emph{Power is not guaranteed under steering.} A planner heading for states where the same error is harder to see delays detection by about $1.5\times$ in our study; we do not bound this in general.
\item \emph{Transients and ramps.} A cumulative e-process detects sustained changes; a $150$-round pulse was detected in $17\%$ of runs and a ramp is detected late.
\item \emph{Scope of the real-video evidence.} The replicated real-video study uses three world-model classes on three shifts (Moving MNIST speed, KTH gait and domain) with 10--30 repetitions; the shifts are large benchmark shifts rather than subtle drifts. Earlier single runs on KITTI driving and UZH-FPV drone video are not replicated and are not used in any claim; subtle drifts and larger models remain untested.
\item \emph{The certificate is an average over logged rounds.} It can be diluted by a short spike inside its window, and it concerns the logged action sequences rather than a new policy's imagined rollouts. No certificate is run on real data here.
\item \emph{Multi-step claims need committed actions.} For $h>1$ the validity argument needs the action sequence fixed at issue time; the endogeneity experiments use $h=1$ and a pivotal statistic.
\item \emph{The gate evidence is one system.} Section~\ref{sec:exp-control} gates a planner on a point mass with three faults, five trained models and 20 episodes per cell. The result depends on the fault suite and on the fallback (zero action is wrong under a persistent disturbance), the gate is slower than a CUSUM gate when a fault is abrupt, and the tolerance is a statistical bound on predictive accuracy, not on task cost. The certified horizon $\hcert_\tau$ was not used in the gate because it is too conservative at the data sizes in an episode, so the gate rests on the revocation frontier alone, with an estimated tolerance.
\end{itemize}

\paragraph{Practical guidance.} Use a statistic whose conditional mean is pivotal if the model provides a calibrated predictive covariance; otherwise add a margin to $\eps_h$ and validate it by replaying held-out in-distribution streams for longer than the deployment horizon. Read the certified horizon over a window rather than the full history. Use the revocation frontier as the alarm
and the certified horizon as the gate; do not treat the first as evidence for the second.

\paragraph{Future work.} Valid pooling of overlapping horizon chains; a tolerance derived from a downstream control-cost budget rather than predictive accuracy; fault-aware fallbacks; windowed e-detectors for transients; and replicated studies of larger neural world models on driving and aerial video.

\paragraph{Reproducibility.} Code, seeds, and the JSON outputs of every table are in the public repository. No new data were collected; the real-video experiments use the public Moving MNIST and KTH datasets.
